% ==============================================================================
%  CHAPTER 1 — INTRODUCTION
% ==============================================================================
\section{Introduction}
\label{sec:intro}

\subsection{Background}

Camera glasses have moved out of research labs and into ordinary shops. Meta's Ray-Ban
glasses look like normal sunglasses but carry a 12-megapixel camera
\cite{gsmarena_display}, and devices such as the Apple Vision Pro and Snap Spectacles
record the world from the wearer's eye level. The market is no longer small. Meta and its
manufacturing partner EssilorLuxottica reported selling about 2 million Ray-Ban Meta
glasses between the launch in October 2023 and February 2025, and about 7 million Ray-Ban
and Oakley smart glasses in 2025 alone, more than three times the number sold in 2024
\cite{uploadvr2026sales}. Each of those devices is a camera that sits on someone's face,
points wherever they look, and can start recording with a tap or a voice command.

Video recorded this way is called \emph{egocentric} or first-person video: the camera
sees what the wearer sees, and moves whenever the wearer turns their head. Large
first-person video collections such as Ego4D \cite{grauman2022ego4d} show how much of
daily life these cameras capture: conversations at home, meals with friends, shopping,
and crowds in the street. That is what makes the devices useful for memory aids,
navigation, accessibility and live translation, and it is why manufacturers are adding
assistants that look at the scene continuously rather than only when the wearer presses a
button. It is also what makes them a privacy problem, because most of the people in that
footage never chose to be recorded.

\subsection{Bystander Privacy as a Public Concern}
\label{sec:concern}

A smartphone held up to film someone is visible, and the person being filmed can at
least see it happening. Camera glasses remove that signal. The camera sits where a pair of
ordinary glasses would, the wearer does not need to raise anything, and the only
indication that recording has started is a small light on the frame. The people most
affected by this are not the wearers, who bought the device and agreed to its terms, but
the people around them, who did neither.

This concern is not new, and it has repeatedly turned into action rather than staying a
theoretical worry. When Google Glass appeared in 2013 and 2014, some bars and restaurants
asked wearers to take the device off, and the Cinema Exhibitors' Association, whose
members run most cinemas in the United Kingdom, asked customers not to wear it in
auditoriums \cite{nbc2014glass}. The device became so associated with covert recording
that its wearers were given a nickname of their own, and it never reached the general
consumer market.

The same concern has followed the current generation. When Facebook launched its first
Ray-Ban camera glasses in Ireland and Italy in 2021, the Irish Data Protection Commission
and the Italian data protection authority issued a joint statement. They pointed out that
a phone or a camera is visible when it records, whereas these glasses have only a very
small indicator light, and that neither Facebook nor Ray-Ban had shown that the light
works as a warning in real conditions. They asked the company to demonstrate that the
light was effective and to run a public information campaign about the product
\cite{dpc2021view}.

Later events showed why regulators doubted the light. In October 2024, two Harvard
students connected a pair of Ray-Ban Meta glasses to a public face search engine and a
language model, and demonstrated a system that could look at a stranger, find their name
and then search the web for details such as their home address, phone number and
relatives \cite{cox2024ixray}. They did not release the code, and said that their aim was
to show how easily existing products could be combined in this way. In early 2026, a BBC
investigation reported hundreds of videos on social media in which men had approached
women in public and filmed them with smart glasses, and then posted the footage as dating
advice. The women interviewed said they had not noticed any light showing that they were
being recorded, several found the videos only after they had been posted and had received
abusive comments, and in one case the person who posted a video asked the woman to pay
for it to be removed \cite{bbc2026glasses}. In May 2026, the Attorney General of Texas
opened an investigation into Meta's camera glasses, stating that the recording light is
easily hidden, that it is not active when the glasses are in a mode that keeps them
processing video for the assistant, and that the devices could collect biometric data and
record people without their knowledge or consent \cite{cbs2026texas}.

Research with bystanders points the same way. In a field study, Denning et al.\ wore a
mock augmented-reality device in cafes over 12 sessions and interviewed 31 people who had
been near it; bystanders were uneasy about being recorded, and their reactions depended on
what was being captured and why \cite{denning2014}. Surveys report that between 65\% and
90\% of bystanders would take defensive action in sensitive settings \cite{mindthegap2026},
and that the protection people expect depends strongly on context \cite{ohagan2023}.
Wearers feel the tension too. In interviews by Bhardwaj et al., people who use camera
glasses every day said they felt responsible for the people around them, but found the
existing privacy indicators ineffective and had no good tool to protect others
\cite{bhardwaj2024}.

\subsection{The Legal Position and the Right to Complain}
\label{sec:legal}

Bystanders are not only uncomfortable; in many places they also have legal rights over
recordings of themselves. Under the European Union's General Data Protection Regulation,
an image from which a person can be identified is personal data, and a person who
believes their data has been processed unlawfully has the right to lodge a complaint with
a data protection authority \cite{gdpr2016}. The regulation does not apply to purely
personal or household activity, but the Court of Justice of the European Union has read
that exemption narrowly. In the \emph{Ryne\v{s}} case, it held that a private camera
which records a public space, even partly, is not purely personal and falls within data
protection law \cite{rynes2014}. A camera worn on the face records public spaces almost
by definition. Sri Lanka's Personal Data Protection Act, No.~9 of 2022, likewise
provides for a Data Protection Authority whose powers include receiving complaints and
conducting inquiries \cite{srilanka_pdpa2022}.

In practice, therefore, a person filmed by a stranger's camera glasses may be able to
complain to a regulator or ask for the footage to be removed, and the cases above show
that people do complain, to journalists, to platforms, to the police and to regulators.
The question this raises for wearers and manufacturers is not whether bystander privacy
matters, but how a device can respect it without becoming useless to the person wearing
it.

\subsection{The Need for Selective Anonymization}
\label{sec:need}

The people who appear in a wearer's video fall into two very different groups. Some are
\emph{interaction partners}: the friend across the table, the shopkeeper the wearer is
paying, the colleague in the meeting. They know they are with the wearer, they can see the
glasses, and they are part of the moment being recorded. Others are \emph{bystanders}:
strangers walking past, people at the next table, a crowd in a market. They have no
relationship with the wearer and no practical way to object.

The two groups need opposite treatment. A recording of a dinner with friends is only
worth keeping if the friends are visible, and an assistant that is meant to help the
wearer in a conversation has to see the person they are talking to. The strangers at the
next table, on the other hand, gain nothing from appearing in the video and are exposed to
every risk described above: being identified, being posted online, and being searched for
later. A protection that treats every face the same way either fails the bystanders or
fails the wearer.

The need also grows with adoption. With millions of camera glasses sold each year
\cite{uploadvr2026sales}, the number of people recorded without consent grows with them,
and a protection that relies on bystanders noticing a light, or on them carrying a device
of their own, will not scale. What is needed is protection that the glasses apply
automatically, that treats strangers and companions differently, and that does not itself
create a new privacy risk such as a database of recognised faces or a stream of video sent
to a server.

\subsection{Why Existing Protections Fall Short}
\label{sec:existing}

\paragraph{The recording light.}
The protection built into today's products is a small light that turns on while the
camera records. It asks bystanders to notice a tiny light, often in daylight, to know
what it means, and then to act on it. The regulators' statement \cite{dpc2021view}, the
accounts collected by the BBC \cite{bbc2026glasses} and the Texas investigation
\cite{cbs2026texas} all report the same weaknesses: the light is hard to see, it can be
covered, and it does not cover every mode in which the camera is used. Even when it
works, it only tells a bystander that recording is happening; it does nothing to protect
them in the recording.

\paragraph{Blurring every face.}
The second approach is to blur every face, as Google Street View does and as Meta's
EgoBlur tool does for recordings from its research glasses \cite{egoblur2023}. On-device
systems such as SafeAR apply the same idea inside the headset \cite{safear2025}. Blanket
blurring protects bystanders, but it also blurs the friend the wearer is talking to. For a
device whose purpose is to capture and assist social life, a recording in which every
companion is blurred has lost most of its value, and wearers have little reason to keep
such a protection switched on.

\paragraph{Selective systems that need extra hardware.}
Research systems have tried to be selective. BystandAR \cite{bystandar} separates
partners from bystanders using the wearer's eye gaze and a directional microphone array,
and BystandARIA \cite{bystandaria} lets bystanders request blurring with a light-emitting
device they carry. Both show that selective protection is possible, but both depend on
hardware that ordinary camera glasses, or ordinary bystanders, do not have. Eye tracking
and microphone arrays add cost, weight and battery drain, and a scheme that only protects
people who carry a special device protects almost nobody in a real street.

\paragraph{Recognising known people.}
A further option would be to recognise the wearer's friends and blur everyone else. This
needs a database of known faces, which is itself sensitive biometric data, and it is
exactly the kind of capability that regulators and the public object to most strongly
\cite{cox2024ixray,cbs2026texas}. Sending the video to a server for processing has the
same problem, since it moves the very faces that should be protected off the device.

\subsection{Problem Statement}
\label{sec:problem}

Putting these points together, I arrive at the problem in four steps.

\paragraph{Step 1: what is used today.}
Camera glasses rely on a recording light, and recorded footage is protected, if at all,
by blurring every face.

\paragraph{Step 2: why that is not enough.}
The light does not protect bystanders in the recording, and blanket blurring removes the
people the wearer is with, so it defeats the purpose of the device. Blurring every face
continuously also costs computation and battery on a device that has little of either.

\paragraph{Step 3: what others have tried.}
Selective systems separate partners from bystanders, but only with eye tracking,
microphone arrays or devices carried by bystanders, none of which ordinary camera glasses
or ordinary bystanders have.

\paragraph{Step 4: what is missing.}
No published method separates partners from bystanders using only the RGB camera that
every pair of camera glasses already has. Doing so is hard for three reasons. It must not
recognise anyone, because a database of known faces is itself a privacy risk. It cannot
use extra sensors. And it must be light enough for the small processors and batteries of
wearable devices, since sending the video elsewhere would expose the very faces it is
meant to protect.

\begin{center}
\fbox{\begin{minipage}{0.9\linewidth}
\textbf{Problem statement.} Camera glasses lack a way to decide, from their single RGB
camera and without recognising anyone, which people in the video are interacting with
the wearer, so that only bystanders are anonymized, using computation that is realistic
for wearable hardware.
\end{minipage}}
\end{center}

\subsection{Importance of This Research}
\label{sec:importance}

Solving this problem matters to several groups at once. For \textbf{bystanders}, it means
that a stranger's face is blurred before the footage is stored, shared or searched, so the
risks described in Section~\ref{sec:concern} are reduced at the source rather than after a
complaint. For \textbf{wearers}, it means they can record and use an assistant in social
situations without exposing the people around them, and without the recording losing the
friends who made the moment worth keeping. For \textbf{manufacturers}, it offers a
protection that works with the camera already in the glasses, that does not need a face
database, and that keeps the video on the device, which is what regulators have asked for.
And for \textbf{society}, it is a step towards wearable cameras whose default behaviour
respects people who never agreed to be recorded, in line with the idea of privacy by
design \cite{cavoukian2009pbd}.

The research is also timely. The devices are already sold in millions, the complaints and
investigations described above happened within the last two years, and the processors in
current glasses already include hardware for running small neural networks
\cite{qcom_ar1}. A method that
fits within that hardware could be deployed as a software update rather than waiting for a
new generation of devices.

\subsection{Research Aim}

The aim of this research is to design, implement and evaluate a privacy pipeline that
blurs the faces of bystanders in first-person video while leaving the faces of the
wearer's interaction partners visible, using only a single RGB camera, without face
recognition, and with computation suitable for wearable-class hardware.

\subsection{Research Questions}
\label{sec:rq}

\begin{description}
  \item[RQ1.] Can a small set of visual cues taken from a single RGB camera (how large a
        face is, where it is in the frame, which way the head is turned, and whether the
        mouth is moving) separate a wearer's interaction partners from bystanders without
        recognising anyone's identity?
  \item[RQ2.] How well does selective anonymization based on these cues protect bystanders
        and keep interaction partners visible, compared with blurring every face?
  \item[RQ3.] How much computation does the pipeline need, and is it feasible on the CPU of
        wearable-class hardware?
\end{description}

\subsection{Research Objectives}

\begin{enumerate}
  \item Build a face detection and tracking pipeline that keeps a stable, anonymous
        identity for each face in first-person video without storing any biometric data.
  \item Design interaction cues from face geometry, head pose and mouth movement, and a
        decision rule that combines them into a partner or bystander label for each face.
  \item Apply blurring only to faces labelled as bystanders, and produce a blanket-blurred
        version of the same video as a baseline.
  \item Assemble a dataset of first-person video and evaluate the system against human
        judgements, measuring bystander protection and partner preservation separately,
        and measuring how much each cue contributes.
  \item Measure the processing time and power draw of the pipeline and estimate the
        performance on a Raspberry Pi~5 and on current camera-glasses hardware.
\end{enumerate}

Table~\ref{tab:rq_map} shows which objectives and which parts of the thesis answer each
research question.

\begin{table}[htbp]
\begin{singlespace}
\centering
\caption{How the research questions map to objectives and to the thesis.}
\label{tab:rq_map}
\renewcommand{\arraystretch}{1.3}
\begin{tabular}{@{}lll@{}}
\toprule
\textbf{Question} & \textbf{Objectives} & \textbf{Answered in} \\
\midrule
RQ1 (cues)        & 1, 2, 4 & Sections~\ref{sec:res_ablation}, \ref{sec:res_errors}, \ref{sec:disc_cues} \\
RQ2 (accuracy)    & 3, 4    & Sections~\ref{sec:res_2a} to \ref{sec:res_baseline}, \ref{sec:disc_tradeoff} \\
RQ3 (computation) & 5       & Sections~\ref{sec:res_speed}, \ref{sec:disc_hw} \\
\bottomrule
\end{tabular}
\end{singlespace}
\end{table}

\subsection{Scope}

I work only with the visual channel. The system does not use audio, depth sensing or eye
tracking, because most consumer glasses do not expose them to third-party software. The
decision rule is a hand-designed heuristic rather than a trained model, since no labelled
dataset of partners and bystanders existed when I started. I evaluate on recorded
first-person video from the Ego4D dataset and from YouTube rather than footage recorded
for this study. Processing speed and power are measured on a laptop CPU and projected to
the Raspberry Pi~5 and to camera-glasses hardware, which I did not have access to. The
thesis is about technical protection; it does not give legal advice, and the legal
position in Section~\ref{sec:legal} is described only to show why the problem matters.

\subsection{Contributions}

\begin{itemize}
  \item A selective anonymization pipeline that separates interaction partners from
        bystanders using only face geometry, head pose and mouth movement from one RGB
        camera, with no face recognition and no stored identity.
  \item An offline second pass for recorded video that decides once per person, keeping
        confirmed partners, and the people who stay with them, visible for their whole time
        on screen, and rejoining people who briefly leave the view without recognising them.
  \item An evaluation protocol for this task: blind labelling of short video clips with
        sound by several annotators, agreement statistics, separate privacy and utility
        measures with confidence intervals, a check of every system decision against the
        rendered output, and an ablation that measures each cue on the same labels.
  \item A breakdown of where processing time and power are spent, with projections to a
        Raspberry Pi~5 and to current camera-glasses hardware.
\end{itemize}

\subsection{Thesis Structure}

The rest of the thesis is organised as follows.

\textbf{Chapter~\ref{sec:lit}, Literature Review,} first gives the background the work
builds on: bystander privacy and how it has been studied, first-person vision, face
detection, object tracking, facial landmarks and the hardware of edge devices and camera
glasses. It then reviews the systems closest to mine, from blanket blurring to selective
systems that use gaze, audio or devices carried by bystanders, and ends by stating the
research gap that this thesis addresses.

\textbf{Chapter~\ref{sec:method}, Methodology,} describes the research design, the
dataset and the tools, and the theory behind each stage of the pipeline. It then explains
the design of the pipeline stage by stage, the decision rule and its thresholds, the
offline second pass for recorded video, and the evaluation: the three labelled sets, the
measures of bystander protection and partner preservation, the agreement and interval
statistics, and the speed, power and projection methods.

\textbf{Chapter~\ref{sec:results}, Results and Discussion,} reports the agreement between
annotators, the accuracy of both versions of the system on each labelled set, the
contribution of each cue, the effect of the area threshold, an analysis of the errors,
the offline second pass, the improvements tested afterwards, the robustness of the
results, and the processing time, power draw and hardware projections. It then discusses
what the results mean, compares them with prior work and with real camera-glasses
hardware, and sets out the limitations of the study.

\textbf{Chapter~\ref{sec:conclusion}, Conclusion,} answers the three research questions,
summarises the contributions and outlines future work.
